\section{Link self-energy contribution and ultraviolet divergences}

The simplest diagram corresponds to the
self-energy correction for the gauge
link (see Fig. \ref{linkself}). Its calculation is the same as in case of the quark bilocal operators
(see, e.g., Ref. \cite{Radyushkin:2017lvu}).
At one loop, one should just the change the color factor $C_F \to C_A$. Thus, we have
\begin{align}
\Gamma_ \Sigma (z) = & ({i} g)^2\,C_A \,\frac12 \, \int_0^1 \dd t_1 \, \int_0^1 \dd t_2 \,
z^\mu z^\nu \, D_{\mu \nu} ^c [ z (t_2-t_1) ] \ ,
\label{self}
\end{align}
where $D_{\mu \nu} ^c (z) =g_{\mu \nu} /4\pi^2 z^2$
is the Feynman-gauge gluon propagator in the coordinate representation.
The resulting integrals over the link parameters $t_1, t_2$
\begin{align}
\int_0^1 \dd t_1 \, \int_0^1 \frac{\dd t_2}{(t_2-t_1)^2} \,
\label{t1t2}
\end{align}
diverge
when $t_1 \sim t_2$, i.e., when
the endpoints $t_1z$ and $t_2 z$ of the gluon propagator are close to each other. So,
one may suspect that this divergence has an ultraviolet origin.
To see that this is the case, we use the dimensional regularization (DR) \cite{tHooft:1972tcz}
in the UV region, switching to
$d$ dimensions.
As a result, the gluon propagator in the coordinate space acquires an extra factor $(-z^2)^{2-d/2}$.
This results in an extra $(t_2-t_1)^{2-d/2}$ factor in Eq. (\ref{t1t2}), and the integral there converges
for sufficiently small $d$.

\begin{figure}[t]
\centerline{\includegraphics[width=2in]{images/linkgluon}}
\vspace{-8mm}
\caption{Self-energy-type correction for the gauge link.
\label{linkself}}
\end{figure}

To preserve gauge invariance, our calculations were made using
massless gluons and the dimensional regularization.
However, in the case of the link self-energy diagram, the use of DR
(which is basically just a mathematical trick) is rather misleading in a couple of points.

The relevant subtleties may be illustrated by using the Polyakov prescription
\mbox{$1/z^2 \to 1/(z^2-a^2)$} for the gluon propagator in the coordinate representation \cite{Polyakov:1980ca}
(see also Refs. \cite{Chen:2016fxx,Radyushkin:2017lvu}).
It
softens the gluon propagator at intervals $-z^2 \lesssim a^2$,
and eliminates its singularity at $z^2=0$. In this respect,
it is similar to the UV regularization produced by a finite lattice spacing, and gives
\begin{align}
\Gamma_ \Sigma (z,a) = & - g^2\,C_A \,\frac{z^2}{8 \pi^2}
\, \int_0^1 \dd t_1 \, \int_0^1 \, \frac{\dd t_2}{ z^2 (t_2-t_1)^2 - a^2} \ .
\label{selfa}
\end{align}
The regularized integral vanishes
on the light cone \mbox{$z^2=0$} and converges for spacelike $z$.
Taking $z=z_3$ and
calculating the integrals gives \cite{Chen:2016fxx,Radyushkin:2017lvu}
\begin{align}
\Gamma_\Sigma (z_3,a) = -\,C_A \,\frac{\alpha_s}{2 \pi}
&
\left [
\,
2 \frac{z_3}{a} \, \tan
^{-1}\left(\frac{z_3}{a}\right) - \ln
\left(1+ \frac{z_3^2}{a^2}\right) \right ] \ .
\label{selfex}
\end{align}
The result contains a linear $\sim 1/a$ divergence that is
missed if one uses the DR. Furthermore,
for a {\it fixed} $a$ and small $z_3$ it behaves
like $z_3^2/a^2$, i.e., $\Gamma_\Sigma (z,a)$ vanishes for $z_3=0$, as expected: there is no link if $z_3=0$.
It also vanishes on the light cone $z^2=0$.

The fact that $\Gamma_\Sigma (z_3=0,a) =0$ means that,
for a fixed $a$,
this term gives no corrections to
the local limit of the $G_{\mu \alpha} (z) \, [z,0]\, G_{\lambda \beta} (0)$
operator, e.g., to
the energy-momentum tensor (EMT).
Since the matrix element
of the EMT gives the fraction of the hadron momentum carried by the gluons,
the link self-energy correction does not change this fraction. This is a natural
phenomenon in the absence of
the gluon-quark transitions.

However, if one formally takes the $a\to 0$ limit for a fixed $z_3$ in Eq. (\ref{selfex}),
then $\ln \left(1+{z_3^2}/{a^2}\right)$ converts into the expression
$\ln \left({z_3^2}/{a^2}\right)$ singular for $z_3=0$.
Similarly, using the DR, one faces an outcome proportional to
\begin{align}
(-z^2\mu_{\rm UV}^2)^{\euv}/\euv= 1/\euv + \ln (-z^2 \mu_{\rm UV}^2) + \ldots \, ,
\end{align}
where $\mu_{\rm UV}$ is the scale accompanying this UV dimensional regularization.
Again, the starting expression vanishes for $z^2=0$, but
renormalizing it by a subtraction of the $1/\euv$ pole,
one may apparently conclude that, in addition to the UV divergence, this diagram
contains a singularity on the light cone $z^2=0$.

For this reason, in our DR results we will explicitly separate
the \mbox{$z^2$-dependence} induced by the UV singular terms
(that actually vanish on the light cone)
and
that present in the DGLAP-evolution
logarithms $\ln (-z^2 \mu_{\rm IR}^2)$, where $\mu_{\rm IR}$ is the scale
associated with the DR regularization of the collinear singularities.

The main difference is that if, instead of DR, one regularizes collinear singularities
by using a physical IR cut-off $\Lambda$ (like nonzero gluon virtuality or gluon mass),
the one-loop result, proportional to the modified Bessel function $K_0 (\sqrt{-z^2 \Lambda^2})$, remains
singular for $z^2=0$, unlike the UV-induced logarithm
$\ln (1-z^2/a^2)$.

In the case of the link self-energy diagram, we have UV singularities only.
Its correction to the $G_{\mu \alpha}(z)G_{\lambda \beta}(0)$ operator is given by
\begin{align}
&
-{g^2N_c \over 4\pi^2[(-z^2 \mu_{\rm UV}^2 +i\epsilon)]^{{d\over 2}-2}} {\Gamma\big({d/ 2}-1\big) \over (3-d)(4-d)}G_{\mu\alpha}(z)G_{\lambda \beta }(0) \ ,
\label{selfAD}
\end{align}
where the $1/(3-d)(4-d)$ factor results
from the integral
\begin{align}
&
\!\int_0^1\! \dd t_1 \!\int_0^{t_1}\! \dd t_2\, (t_1-t_2)^{2-d} ={1 \over (3-d)(4-d)}
\nonumber
\end{align}
produced by the DR of the gluon propagator \mbox{$D^c (t_1z - t_2 z)$.}
The pole for $d=3$ ($d=4$) corresponds to the linear (logarithmic) UV divergence
in Eq. (\ref{selfex}).
